\begin{abstract}

Las imágenes de Radar de Apertura Sintética (SAR) son fundamentales para la observación de la Tierra debido a su capacidad de adquisición independiente de las condiciones climáticas y de iluminación. Sin embargo, estas imágenes se ven afectadas por el ruido \emph{speckle}, un patrón granular multiplicativo que degrada su calidad, y dificulta su interpretación.

Este trabajo presenta un enfoque novedoso para la eliminación de ruido \emph{speckle} mediante el uso de autoencoders convolucionales. La principal innovación radica en el uso de la distribución $\mathcal{G}_I^0$ para generar sintéticamente pares de imágenes ruidosas y limpias destinadas al entrenamiento supervisado, modelando fielmente las características estadísticas del ruido \emph{speckle} real. Complementariamente, desarrollamos un sistema flexible basado en archivos de configuración que facilita la experimentación con diferentes arquitecturas.

Los experimentos revelan que las arquitecturas más profundas obtienen consistentemente mejores resultados sin evidencia de \emph{overfitting}. Los modelos entrenados con datos sintéticos mostraron buen desempeño en imágenes de test sintéticas, pero limitaciones en la generalización a imágenes SAR reales. Para abordar esto, exploramos el entrenamiento con pares de imágenes reales obtenidas mediante promediado temporal, logrando mejoras significativas y desempeño comparable a otros ya conocidos.

Este trabajo establece las bases para el uso de autoencoders en el filtrado de \emph{speckle} y demuestra tanto el potencial como las limitaciones de los datos sintéticos generados mediante modelos estadísticos.

\vspace{0.5cm}

\noindent\textbf{Palabras clave:} Radar de Apertura Sintética (SAR), ruido \emph{speckle}, \emph{despeckling}, autoencoders, redes neuronales convolucionales, distribución $\mathcal{G}_I^0$, procesamiento de imágenes, filtrado de ruido.

\end{abstract}